As written:
For each $n\in\mathbb{N}$, define $g(n)$ to be the minimal integer such that for every set $A\subset\mathbb{N}\cap[2,\infty)$ with $|A|=n$, and every set $I$ of $\max(A)$ consecutive integers, there exists $B\subseteq I$ with $|B|=g(n)$ such that
\[
\prod_{a\in A} a\ \big|\ \prod_{b\in B} b.
\]
Question: is $g(n)\le (2+o(1))n$? perhaps even $g(n)\le 2n$?

\textbf{Ambiguity/misstatement.}
The ``$|B|=g(n)$'' clause is incompatible with the subsequent discussion in the problem text.
Indeed, because $B\subseteq I$ we always have $|B|\le |I|=\max(A)$.
But for fixed $n$, we may choose $A=\{2,3,\dots,n+1\}$, for which $\max(A)=n+1$.
If the definition required existence of a subset $B$ of \emph{exactly} size $g(n)$ for \emph{every} such $A$, then necessarily $g(n)\le n+1$, which would make an asymptotic bound like $g(n)\le (2+o(1))n$ vacuous.

\textbf{Minimal corrected statement (standard convention).}
Define $g(n)$ as the minimal integer such that for every $A$ and $I$ as above, there exists $B\subseteq I$ with
\[
|B|\le g(n)\quad\text{and}\quad \prod_{a\in A} a\mid \prod_{b\in B} b.
\]
This ``$\le g(n)$'' formulation is monotone: if a set $B$ works, any superset also works.
I interpret the problem under this corrected definition.

2)
